\section{Background and related work}
\label{sec:related}

\paragraph{The format.} PDF is specified by ISO 32000-1~\cite{iso32000}. A file is a
sequence of indirect objects, a page tree under a catalogue, and a cross-reference table
that maps object numbers to byte offsets. Two properties matter here. First, the format
admits \emph{incremental updates}: a new section may be appended at the end of the file with
its own cross-reference table and a \texttt{/Prev} pointer to the previous one, so the
original bytes never move. Second, page attributes are \emph{inheritable} along the page
tree: a page that omits \texttt{/MediaBox} takes the nearest ancestor's.

Streams are stored as a dictionary followed by raw bytes whose length the dictionary
declares. The bytes may be compressed by any of several filters, and objects may be packed
into a compressed container called an object stream. A writer that wants to copy a stream
without decoding it must therefore preserve the bytes exactly and copy the filter
description; a writer that has to rewrite a reference \emph{inside} an object stream cannot
do so without decoding and re-encoding the container.

\paragraph{Rendering engines.} pdf.js~\cite{pdfjs} is a JavaScript PDF renderer, Apache-2.0
licensed, and the rendering engine of Firefox. It exposes a page model and a canvas
rasteriser, plus a text-extraction API that yields positioned glyph runs. MuPDF~\cite{mupdf}
is a C engine, dual-licensed AGPL/commercial, with annotation support. PDFium is the engine
inside Chromium. These are renderers first; writing is either absent or a secondary concern.

\paragraph{Readers and annotators.} Okular~\cite{okular} (KDE, Poppler-backed) and Evince
(GNOME) are mature open-source readers with annotation support that write annotations into
the document through their respective back ends. Apache PDFBox~\cite{pdfbox} and iText are
libraries with full writing support, used by other applications rather than being
applications themselves. PDF-XChange Viewer~\cite{pdfxchange}, whose feature set this project
follows, is a commercial Windows viewer that is no longer developed.

The gap this project occupies is deliberately narrow: an \emph{application}, in which the
reader, the writers and the annotation model are open and readable together, with no engine
dependency beyond rasterisation, and with the writing logic documented as a set of invariants
rather than as behaviour to be discovered. We do not claim to be better at rendering, and we
do not compare rendering quality against the systems above; \S\ref{sec:eval} explains why
such a comparison would be neither fair nor useful at this stage.

\paragraph{Why write at all.} A sidecar annotation file is simpler in every engineering
respect: no incremental updates, no object numbering, no risk of corrupting the document. It
is also invisible to every other application. The design decision recorded here is that the
document is the artefact that travels, so the annotations must travel inside it. The price is
paid in \S\ref{sec:pdf-layer}.

\begin{table}[t]
  \centering
  \caption{Dependencies. The application's only runtime dependency is the rendering engine.
  Everything else is a build or verification tool and ships in no artefact.}
  \label{tab:deps}
  \small
  \begin{tabular}{@{}llll@{}}
    \toprule
    Package & Role & Version & Licence \\
    \midrule
    \texttt{pdfjs-dist} & rendering and text extraction & 5.x & Apache-2.0 \\
    \midrule
    \multicolumn{4}{@{}l}{\emph{build and verification only}} \\
    \texttt{electron} & desktop host & 44.x & MIT \\
    \texttt{electron-builder} & packaging & 26.x & MIT \\
    \texttt{vite} & bundler & 7.x & MIT \\
    \texttt{typescript} & type checker (erased at build) & 5.9 & Apache-2.0 \\
    \texttt{playwright} & browser and Electron driving & 1.64 & Apache-2.0 \\
    \bottomrule
  \end{tabular}
\end{table}
